%----------------------------------------------------------------
% Resumo e Abstract
%----------------------------------------------------------------
\begin{resumo}{estudo autônomo, avaliação formativa, geração de questões, modelos de linguagem, visão computacional no navegador, privacidade}
O estudo autônomo diante de uma tela reúne dois problemas: a atenção se dispersa sem que o estudante perceba, e não há quem verifique, no momento da leitura, se o conteúdo foi compreendido. Modelos de linguagem permitem gerar perguntas sobre qualquer material, e a visão computacional permite estimar, pela câmera, se o estudante olha para a tela. As duas soluções, porém, trazem riscos pouco medidos: perguntas que o material não sustenta, e monitoramento que envia imagens do rosto a servidores. Este trabalho apresenta o \ferramenta, uma ferramenta de apoio ao estudo autônomo em que o próprio estudante envia seus PDFs, recebe questões de múltipla escolha geradas a partir de fragmentos específicos desse material e acompanha, em um painel de automonitoramento, indicadores de engajamento visual calculados inteiramente no navegador, sem que nenhuma imagem deixe o dispositivo. A pergunta que orienta o trabalho é se essa combinação é viável em computadores comuns de estudantes e com que qualidade, desempenho e garantia de privacidade ela funciona. Seguindo a \estr{Design Science Research}, o artefato é avaliado por um protocolo técnico reproduzível, sem participantes humanos: inspeção cega das questões geradas, comparadas a duas linhas de base; medição de desempenho em hardware heterogêneo a partir de um vídeo rotulado; validação do indicador de engajamento contra estados comportamentais roteirizados; e verificação automatizada do tráfego de rede. \resultado{síntese de uma a duas frases com os principais números: proporção de questões aceitáveis e IC, desempenho mediano no navegador, acurácia balanceada do indicador, resultado do teste de privacidade}. Os resultados delimitam em que condições a abordagem é viável e onde ela falha, e oferecem um protocolo reutilizável para avaliar ferramentas educacionais que combinam modelos de linguagem e sensoriamento no cliente.
\end{resumo}

\begin{abstract}{self-directed study, formative assessment, question generation, language models, in-browser computer vision, privacy}
Self-directed study in front of a screen combines two problems: attention drifts without the student noticing, and nobody checks, while reading, whether the content was understood. Language models can generate questions about any material, and computer vision can estimate, through the camera, whether the student is looking at the screen. Both solutions, however, carry poorly measured risks: questions the material does not support, and monitoring that sends face images to servers. This work presents \ferramenta, a self-directed study tool in which students upload their own PDFs, receive multiple-choice questions generated from specific fragments of that material, and follow, on a self-monitoring dashboard, visual engagement indicators computed entirely in the browser, with no image ever leaving the device. The guiding question is whether this combination is feasible on ordinary student computers, and with what quality, performance and privacy guarantees. Following Design Science Research, the artefact is evaluated through a reproducible technical protocol with no human participants: blind inspection of generated questions against two baselines; performance measurement on heterogeneous hardware using a labelled video; validation of the engagement indicator against scripted behavioural states; and automated verification of network traffic. \resultado{one-to-two-sentence summary of the main figures}. The results delimit when the approach is feasible and where it fails, and provide a reusable protocol for evaluating educational tools that combine language models with client-side sensing.
\end{abstract}
